% LaTeX template for Summer Schools
%
% !!!!!!!!!!!!   INSTRUCTIONS  !!!!!!!!!!!!!!

% 0) Submit your work in tex format (not pdf). It will be pasted into the proceedings file.
%
% 1) please name your file yourlastname.tex, e.g. thiele.tex.
%
% As several files will be concatenated, please 
% use some discipline as to the following:
%
% 2) Whenever you use \label{} to get automatic cross references
% through \ref{} (this is the preferred option for cross references)
% please add your initials to the label such as 
% \label{SLEct} for the Stochastic Loewner equation
% with initials of author Christoph Thiele
% Same with other citations such as in bibitem.
%
% 3) Please STRONGLY avoid using \def or \newcommand unless really necessary.
% We do have macros for black board bold. If you use your favorite
% macros while preparing your summary, please expand them
% (replace by the original definition) everywhere in your file.
% This will save me the work of doing the very same thing.
% Thank you! If you have to use \def, please also add your
% initials to the definition.
%
% 4) if you want to compile the header of the document, 
% uncomment remove the corresponding paragraph signs below
%
% 5) There is a sample lecture below. For the header 
% it is best to keep most of the
% commands and just change the name, title, text. etc
%
% 6) Please follow the conventions below in terms of capitalization of 
% headings etc:
% Only the beginning of a sentence and names are capitalized.


\documentclass[12pt]{article}
\usepackage{amssymb,amsmath,amsthm}
\usepackage{esint}
\usepackage{mathrsfs}
\usepackage{mathtools}
\usepackage{comment}
\usepackage{bbm,dsfont}

\newtheorem{theorem}{Theorem}
\newtheorem{lemma}[theorem]{Lemma}
\newtheorem{corollary}[theorem]{Corollary}
\newtheorem{conjecture}[theorem]{Conjecture}
\newtheorem{definition}[theorem]{Definition}
\newtheorem{example}[theorem]{Example}
\newtheorem{remark}[theorem]{Remark}
\newtheorem{proposition}[theorem]{Proposition}

\newcommand{\talktitle}[1]{\section{#1}}
\newcommand{\talkafter}[1]{\textbf{After #1} \addcontentsline{toc}{subsection}{after #1}}
\newcommand{\talkspeaker}[2]{\begin{center}
\textit{A summary written by #1}
\end{center}
\addcontentsline{toc}{subsection}{#1, #2}
}

\usepackage{graphicx}
\usepackage[utf8]{inputenc}
\usepackage[T1]{fontenc}
\usepackage[hidelinks]{hyperref}

\newcommand*{\Z}{\mathbb{Z}}
\newcommand*{\Q}{\mathbb{Q}}
\def\C{\mathbb{C}}
\newcommand*{\N}{\mathbb{N}}
\newcommand*{\R}{\mathbb{R}}

% Absolute values and norms using mathtools.
% \\[lr][vV]ert produces correct spacing, as opposed to | and \|.
\DeclarePairedDelimiter\abs{\lvert}{\rvert}
\DeclarePairedDelimiter\norm{\lVert}{\rVert}

\title{Nodal Domains and Landscape Functions}

% \author{Summer School, Kopp}
% \thanks{supported by Hausdorff Center for Mathematics, Bonn}

% \date{October 2022}



\begin{document}

% \maketitle

% {
% \center{Organizers:}
% \center{
% Christoph Thiele, Universit\"at Bonn}
% \center{
% Olli Saari, Universit\"at Bonn}
% \center{$\ $}
% }
% \newpage
% \tableofcontents

% \newpage


\talktitle{Gaps and bands for Schr\"odinger operators}
\talkafter{J. Garnett and E. Trubowitz \cite{gtkc}}
\talkspeaker{Kaihua Cai}{Caltech}

\setcounter{equation}{0}
\setcounter{theorem}{0}

\begin{abstract}
{ We give a simple necessary and sufficient condition on
the length of the gaps corresponding to one dimensional periodic
Schr\"odinger operators. }
\end{abstract}

Macros demonstration:
\[
\R,\Z,\C,\Q,\N,\abs{x},\abs*{\sum x},\norm{f}, \norm*{\sum f}.
\]

\subsection{Introduction}

Let $q(x)$ be the periodic extension to the whole line of a real
valued function in $L^2_R [0,1]$. The spectrum of $- (d^2/dx^2)
+q(x)$ acting on $L^2(R)$ is the union of purely absolutely
continuous bands $B_n(q), n \geq 1$. It is well known that the
bands may touch but never overlap. We assume that the bottom of
the first band is at $0$. The complement of the spectral band in
$(0,+\infty) $ is a sequence of open intervals, called the gaps.
To each set of bands $B_n(q), n \geq 1$, we associate the sequence
of nonnegative numbers

$$ \gamma_1(q), \gamma_2(q),\cdots $$

\noindent where $  \gamma_n(q)$ is the distance between the top of
the n-th band and the bottom of the next. The main theorems of
this paper is to describe the set of all possible configurations
of bands by understanding the distribution of gaps.



\begin{theorem}

Let $\gamma_n, n \geq 1,$ be any nonnegative sequence satisfying $
\Sigma _{n \geq 1 } \gamma_n^2  < \infty  $. Then, there is a way
of placing the sequence of open tiles of lengths $\gamma_n$, in
order on the positive axis $ (0, \infty)$ so that the compliment
is the set of bands for a function $q$ in $L^2_R [0,1]$. In other
words, the map

$$q \rightarrow \gamma (q)= (\gamma_n(q), n \geq 1), $$

\noindent  from $L^2_{\R} [0,1]$ to $(l^2)^+ $, is onto.
\end{theorem}

It is natural to ask how many different ways a sequence of open
tiles, whose lengths are $\gamma_n, n \geq 1$, can be placed so
that the complement is a set of bands. This is answered by the
following theorem:

\begin{theorem}

There is just one way to place a sequence of open tiles,
satisfying the hypothesis of Theorem 1, on the positive real axis
so that they are genuine gaps.

\end{theorem}

To prove these two theorems, we use a characterization of bands
due to Marcenko and Ostrovskii. They identify band configurations
with slit quarter planes.

Let $\mu_n (q), n \geq 1, $ and $\nu_n(q), n \geq 0,$ be the
Dirichlet  and Neumann spectrum of $q$ in $L_R^2[0,1]$, that is,
the spectra of

\begin{equation} -y''+q(x)y= \lambda y  \label{eq:mainkc}
\end{equation}

\noindent for the boundary condition $ y(0)=0, y(1)=0 $  and $
y'(0)=0, y'(1)=0 $ respectively. If $q$ is an even function, then
$ \gamma_n (q)= |\mu_n (q)-\nu_n(q) |$. We define the signed gap
lengths of $q$ in $E_0$, the subspace of even functions in
$L_R^2[0,1]$ with mean $0$, to be the sequence $(\mu_n
(q)-\nu_n(q)), n \geq 1$. We have the following:

\begin{theorem}
The map from $q$ to its signed gap lengths is a real analytic
isomorphism between $E_0$ and $l^2$.

\end{theorem}

In fact, we obtain three real analytic isomorphisms between the
three spaces $E_0$,  $l^2$ and $l^2_1$, the space of real
sequences $ \{h_n \} $, satisfying $ \Sigma n^2 h_n ^2 < \infty $.

\subsection{Preliminaries}

\subsubsection{More levels}

Some people need more levels, which can be generated with subsubsection

\subsubsection{Here we go} 

Let $y_1(x,\lambda, q )$ and $y_2(x,\lambda, q )$ be the solutions
of \eqref{eq:mainkc} satisfying

$$y_1(0,\lambda )=y'_2(0,\lambda )=1, \quad \quad y'_1(0,\lambda )=y_2(0,\lambda )=0 $$

\begin{thebibliography}{03}

\bibitem[GT]{gtkc} Garnett, J. and Trubowitz, E.,
\emph{ Gaps and bands of one dimensional
periodic Schr\"odinger operators.} 
Comment. Math. Helv. {\bf 62} (1987), no.~1, 18--37;


\bibitem[MO]{MOkc} Marcenko, V. A. and Ostrovshii, I. V.,
\emph{ A Characterization of the spectrum of Hill's operator.}
Math. USSR-Sbornik, 97 (1975), pp. 493-554.




\end{thebibliography}

\noindent \textsc{Kaihua Cai, Caltech}\\
\textit{email:} \texttt{cai@caltech.edu}


%\newpage

\end{document}
